% Related systems and baselines.
% This is a paper fragment and assumes tabularx, booktabs, hyperref, and natbib.

\begin{table*}[t]
\centering
\small
\setlength{\tabcolsep}{4pt}
\renewcommand{\arraystretch}{1.08}

\begin{tabularx}{\textwidth}{
    @{}
    >{\raggedright\arraybackslash}X
    >{\raggedright\arraybackslash}p{0.135\textwidth}
    c
    >{\raggedright\arraybackslash}p{0.16\textwidth}
    >{\raggedright\arraybackslash}p{0.145\textwidth}
    c
    @{}
}
\toprule
\textbf{System}
& \textbf{Primary focus}
& \textbf{Decl.}
& \textbf{Targeting}
& \textbf{Validation}
& \textbf{Export} \\
\midrule
PyTorch/\texttt{safetensors} scripts
& Scripts
& --
& Manual
& Manual
& Manual \\

MergeKit~\citep{goddard2024mergekit}
& Merging/MoE
& Yes
& Layer slices/filters
& Limited
& \textbf{Yes} \\

LoRA/adapter tooling~\citep{hu2021lora}
& Adapters
& Partial
& Adapter-specific
& Limited
& Partial \\

\texttt{safetensors}/HF utilities
& Storage/loading
& --
& --
& --
& \textbf{Yes} \\

\href{https://docs.pytorch.org/docs/stable/distributed.checkpoint.html}{PyTorch DCP}
& Distributed I/O
& --
& State dict/planner
& Metadata/load
& \textbf{Yes} \\

\href{https://github.com/google/orbax}{Orbax}
& JAX checkpointing
& Partial
& PyTree/partial
& Restore/metadata
& \textbf{Yes} \\

\href{https://pypi.org/project/torch-state-bridge/}{\texttt{torch-state-bridge}}
& Key rewriting
& Partial
& Rules/regex
& Preview/diff/collision
& Manual \\

\href{https://dnhkng.github.io/posts/rys/}{RYS}/\href{https://github.com/alainnothere/llm-circuit-finder}{LLM Circuit Finder}
& Layer-path surgery
& Partial
& Layer indices
& Task probes
& \textbf{Yes} \\

TransformerLens~\citep{nanda2022transformerlens}
& Interpretability
& --
& Model-internal
& Analysis
& -- \\

nnsight~\citep{fiottokaufman2024nnsight}
& Tracing
& --
& Activation
& Analysis
& -- \\

\midrule
\textbf{BrainSurgery}
& \textbf{Rewriting}
& \textbf{Yes}
& \textbf{Regex/structured}
& \textbf{Assert/diff}
& \textbf{Yes} \\
\bottomrule
\end{tabularx}

\caption{Qualitative comparison with related tooling. MergeKit, Orbax,
\texttt{torch-state-bridge}, PyTorch Distributed Checkpoint (DCP), and
RYS/LLM Circuit Finder overlap with particular parts of the workflow, but differ
in scope: merging, JAX restoration, key rewriting, distributed checkpoint I/O,
or layer-path surgery, respectively. BrainSurgery provides general-purpose
checkpoint rewriting through declarative plans with structured targeting,
executable postconditions, checkpoint diffs, and export.}
\label{tab:tool-comparison}
\end{table*}

\paragraph{Overlapping systems.}
MergeKit is a substantive baseline rather than only a generic merging utility:
its YAML configurations support layer slicing and tensor filters, and its broader
tooling includes out-of-core execution, LoRA extraction, MoE construction, and
multi-stage workflows. Its abstraction nevertheless remains centered on merging
and related model-building tasks. \texttt{torch-state-bridge} is narrower but
more directly overlaps key transformation: it offers rule and regex-capture
mappings, arithmetic substitutions, composable pipelines, previews, diffs, and
collision detection over in-memory PyTorch state dictionaries. Because its
currently published artifact is an early software release, comparisons should pin
the exact package version.

Orbax provides the closest JAX-oriented checkpointing comparison. Its current V1
API supports partial PyTree restoration and restore-time dtype, shape, and
sharding changes. The legacy regex/value-function transformations API is
deprecated, while broader arbitrary V1 model surgery is documented as planned;
claims should therefore identify the compared Orbax API version. PyTorch DCP is
primarily an I/O and scaling baseline, providing parallel multi-rank checkpoint
save/load and load-time resharding rather than a general tensor-rewrite language.
Although BrainSurgery supports DCP directories, its present adapter uses
single-process operations over plain tensor state dictionaries, so this support
does not establish parity for multi-rank execution, optimizer states, or
resharding.

Finally, \href{https://dnhkng.github.io/posts/rys/}{RYS} and
\href{https://github.com/alainnothere/llm-circuit-finder}{LLM Circuit Finder}
illustrate a concrete structural-surgery use case: they alter layer execution
paths and can materialize duplicated layers in GGUF checkpoints before downstream
evaluation. They are best treated as motivating related software rather than a
general framework baseline, and their current blog/repository status warrants
separating capability comparisons from their unreviewed empirical claims.
